
Der Prozess-Fit beschreibt die Passung zwischen den Prozessanforderungen eines ERP-Systems und den vorhandenen 
Geschäftsprozessen eines KMU. Entscheidend ist, ob zentrale Abläufe ausreichend klar beschrieben, nachvollziehbar 
und im System abbildbar sind. 

Damit ihre Integrationswirkung realisiert werden kann, müssen Abläufe zwischen verschiedenen Funktionsbereichen 
aufeinander abgestimmt sein. Dies betrifft beispielsweise Beschaffung, Produktion, Lagerhaltung, Vertrieb, 
Rechnungswesen oder Qualitätsmanagement \parencite{Alaskari2021,Leyh2015}. Sind Prozesse nur informell geregelt, 
stark personenabhängig oder zwischen Abteilungen uneinheitlich, steigt das Risiko, dass sie im ERP-System nur schwer 
abgebildet werden können \parencite{Jha2008}.

Eine wichtige Voraussetzung für einen hohen Prozess-Fit ist daher eine systematische Analyse der bestehenden Abläufe. 
Vor der ERP-Einführung müssen Unternehmen verstehen, welche Prozesse existieren, welche Prozessschritte notwendig sind 
und welche Abhängigkeiten zwischen Abteilungen bestehen. Eine unzureichende Ist-Analyse kann dazu führen, dass Anforderungen 
unvollständig oder falsch definiert werden. Dadurch steigt das Risiko von Fehlkonfigurationen, späteren Anpassungen oder 
ineffizienten Arbeitsabläufen im neuen System \parencite{Leyh2015,Leyh2019}.

Gerade in KMU sind Prozesse häufig weniger formalisiert als in Großunternehmen. Entscheidungen werden oft kurzfristig 
getroffen, Zuständigkeiten sind teilweise informell geregelt und Abläufe beruhen stark auf Erfahrungswissen einzelner 
Mitarbeitender \parencite{Barbieri2025}. Diese Flexibilität kann im Tagesgeschäft vorteilhaft sein, erschwert jedoch die 
ERP-Einführung, da klare Prozesslogiken, definierte Rollen und ein Mindestmaß an Prozessstabilität erforderlich sind 
\parencite{Alaskari2021,Leyh2015}.

Ein weiteres Bewertungskriterium ist die Standardisierbarkeit der Prozesse. ERP-Systeme entfalten ihren Nutzen besonders 
dann, wenn Unternehmen bereit und in der Lage sind, ihre Abläufe zumindest teilweise an bewährte Standardprozesse des 
Systems anzupassen. Werden bestehende Sonderprozesse dagegen vollständig im System nachgebildet, kann dies zu umfangreichen 
Individualanpassungen führen. Diese erhöhen Implementierungsaufwand und Kosten und erschweren häufig spätere Updates, 
Wartung und Systemerweiterungen \parencite{Leyh2015,Barbieri2025}.

Prozess-Fit setzt zudem klare Verantwortlichkeiten voraus. Da ERP-Systeme verschiedene Unternehmensbereiche verbinden, 
können Fehler oder Verzögerungen in einem Prozessschritt Auswirkungen auf nachgelagerte Bereiche haben. Daher muss eindeutig 
festgelegt sein, wer für Prozessschritte, Freigaben, Datenpflege oder Prozessänderungen verantwortlich ist. Fehlen solche 
Zuständigkeiten, kann die Nutzung des ERP-Systems uneinheitlich erfolgen und die angestrebte Transparenz eingeschränkt werden 
\parencite{Leyh2015,Leyh2019,Jha2008}.

Ein hoher Prozess-Fit liegt somit vor, wenn zentrale Abläufe strukturiert, dokumentiert, stabil und weitgehend systemfähig sind. 
Wesentliche Kriterien sind dabei die Dokumentation zentraler Prozesse, klare Rollen und Verantwortlichkeiten, die 
Standardisierbarkeit der Abläufe, die bereichsübergreifende Abstimmung sowie der erwartete Anpassungsbedarf am ERP-System.

